% CHECKED DEDUCTION: full self-reading, archive independent audit and root review passed.
\subsection{Integral odd generation at every Vandiver prime}
\label{sec:integral-classification}

The kernel theorem gives a further integral consequence without locating
the first exceptional weight.  It distinguishes every irregular Vandiver
prime, rather than only the two explicit examples above.

\begin{theorem}[K17-INTEGRAL-CLASSIFICATION]
\label{thm:integral-classification}
Let $p$ be an odd prime satisfying Vandiver, and let
$\mathfrak g_p$ be the actual integral tripod Lie algebra with its
arithmetic weight grading.
If $p$ is regular, $\mathfrak g_p$ is the free $\mathbb Z_p$-Lie
algebra on one generator in every odd weight $m\ge3$.
If $p$ is irregular, there are infinitely many even weights $n$ for which
\[
 \bigl(\mathfrak g_p/[\mathfrak g_p,\mathfrak g_p]\bigr)_n
\]
is a nonzero finite $p$-group.  In particular, no collection of
homogeneous odd integral classes generates $\mathfrak g_p$ as a
$\mathbb Z_p$-Lie algebra.
\end{theorem}

\begin{proof}
The regular-prime assertion is Sharifi's Theorem~1.1 combined with
Brown's rational Deligne--Ihara theorem
\cite[Theorem~1.1, p.~2]{Sharifi2002}\cite[\S\S2.1--2.3]{Brown2019}.
For the irregular case, put
\[
 K=\mathbb Q(\mu_{p^\infty}),\qquad
 H=\operatorname{Gal}(\Omega_p/K),\qquad
 M=H/\overline{[H,H]}.
\]
Theorem~\ref{thm:vandiver-kernel} identifies $H$ with the actual
cyclotomic-kernel image.  Its weight filtration $F^nH$ is therefore
separated.  Proposition~\ref{prop:outer-sat} identifies it with the
actual integral outer filtration.  Write $c$ for complex conjugation.

We first justify that $M^+$ is a nonzero finite free
$\mathbb Z_p$-module.  The Iwasawa and Ferrero--Washington inputs used
in the proof of Theorem~\ref{thm:pole-general} make $H$ free pro-$p$;
thus $M$ has no $\mathbb Z_p$-torsion.  The Kummer comparison
\[
 M^+\simeq (X_\infty^-)^\iota(1)
 \quad\hbox{up to finite kernel and cokernel}
\]
and $\mu(X_\infty^-)=0$ make $M^+$ finitely generated over
$\mathbb Z_p$; hence it is finite free
\cite[Theorem~2]{Iwasawa1981}
\cite[Proposition~3.4.7 and Corollary~3.4.8]{SharifiIwasawaNotes}.

For nonvanishing, let $A_0$ be the $p$-part of the class group of
$\mathbb Q(\mu_p)$ and let $A_\infty$ be the direct limit of the
cyclotomic-layer $p$-class groups under extension of ideals.
Irregularity makes $A_0$ nonzero, and Vandiver makes $A_0^+=0$,
so $A_0^-\ne0$.  The minus-part maps into this direct limit are
injective, and Kummer duality gives an actual isomorphism
\[
 A_0^-\hookrightarrow A_\infty^-\ne0,
 \qquad M^+\cong(A_\infty^-)^\vee(1)
\]
\cite[Propositions~3.4.1 and~3.4.7]{SharifiIwasawaNotes}.
Here $\vee$ is Pontryagin duality; a nonzero discrete torsion
$p$-group has nonzero Pontryagin dual.  Thus $M^+\ne0$, and its
finite free rank is positive.

Let $\pi:H\to M$ be the quotient map and give $M$ the image
filtration
\[
 M^n=\pi(F^nH).
\]
Since $F^1H=H$, this starts with $M^1=M$.
These are closed $\mathbb Z_p$-submodules, because $F^nH$ is
compact.  This filtration is separated.  Indeed, if $x$ belonged
to every $M^n$, the nonempty compact sets
\[
 F^nH\cap\pi^{-1}(x)
\]
would be nested and would have a common point.  That point would
belong to $\bigcap_nF^nH=\{1\}$, forcing $x=0$.
This is the compactness argument needed here; separatedness is
not being assumed to pass automatically to an arbitrary quotient.

There is a surjective graded Lie map to an abelian graded module,
and hence surjections in every degree
\[
 \mathfrak g_p\longrightarrow\operatorname{gr}M,
 \qquad
 (\mathfrak g_p^{\mathrm{ab}})_n
       \twoheadrightarrow M^n/M^{n+1}.                       \tag{IC.1}
\]
Surjectivity follows from the definition of $M^n$; all brackets
map to zero since $\pi$ kills the closed group commutator subgroup.
Complex conjugation acts on $\mathfrak g_{p,n}$ by $(-1)^n$,
so it acts by the same scalar on the quotient in (IC.1).
Since the idempotents $(1\pm c)/2$ are integral, the induced
filtration $M^{n,+}=M^n\cap M^+$ has
\[
 M^{n,+}/M^{n+1,+}=0\quad\text{for odd }n.                   \tag{IC.2}
\]

Each $\mathfrak g_{p,n}$ is a finite free $\mathbb Z_p$-module:
it embeds in the finite free outer Johnson module described in
Proposition~\ref{prop:outer-sat}.  In a fixed weight the bracket
submodule is a finite sum of images of tensor products of finite
free modules, so it is finitely generated and closed.  Thus
$(\mathfrak g_p^{\mathrm{ab}})_n$ is finitely generated.
After tensoring with $\mathbb Q_p$, Brown's theorem makes the
Lie algebra free on odd generators and gives
\[
 (\mathfrak g_p^{\mathrm{ab}})_n\otimes\mathbb Q_p=0
 \quad\text{for even }n.
\]
Exactness of localization and the finite sum defining brackets
justify commuting this abelianization with tensoring.
Therefore every even $(\mathfrak g_p^{\mathrm{ab}})_n$ is a
finite $p$-group, and so is every even quotient
$M^{n,+}/M^{n+1,+}$.

There must be infinitely many nonzero such even quotients.
Otherwise (IC.2) would make the filtration on $M^+$ stationary
from some point onward.  Its intersection is zero, so the
stationary submodule would be zero.  Starting from $M^+$ and
using the finitely many preceding finite quotients would then
make $M^+$ finite, contradicting its positive free rank.
By (IC.1), each of those infinitely many even weights has
nonzero $(\mathfrak g_p^{\mathrm{ab}})_n$.

Finally, a graded Lie algebra generated by homogeneous odd
classes has abelianization generated by their images, all in odd
weights.  Its even abelianization would be zero.  The preceding
nonvanishing therefore excludes every such integral generating
collection, proving the theorem.
\end{proof}

Together with the separately proved all-weight theorem at $2$,
this classifies the existence of free homogeneous odd integral
generators among the primes covered by the Vandiver hypothesis:
they exist at $2$ and at the odd regular primes, and fail at every
irregular Vandiver prime.  The negative assertion does not locate
the first defect weight or its order.  The explicit weight-$32$
example at $37$ and weight-$12$ example at $691$ retain that
additional information.  Kernel-field equality and primitive
full-domain Soul\'e characters can hold at all these irregular
Vandiver primes despite the infinitely many even integral
indecomposable quotients.
